% META: {"id": "TSPE-COMBI-ALG-015", "chapitre": "TSPE-COMBINATOIRE", "type_objet": "algorithme", "capacites_codes": ["C2", "C4"], "programme_alignment": "OFFICIAL_ALGORITHM_EXAMPLE", "mode_creation": "ex_nihilo", "sources_inspiration": [], "status": "generated"}

\section{Algorithmique : engendrer des coefficients binomiaux et des permutations}

% ============================================================
% STRATE 1 --- Probleme, modele, algorithmes
% ============================================================

\margeAppui{
Le programme cite parmi ses \textbf{exemples d'algorithme} : la generation de
la liste des $\binom{n}{k}$ par la relation de Pascal, la generation des
permutations d'un ensemble fini ou le tirage aleatoire d'une permutation, et
la generation des parties a $2$ ou $3$ elements. Il les cite, il ne les impose
pas.
}

\textbf{Le probleme.} Denombrer, c'est compter sans enumerer. Mais pour
verifier un denombrement -- et pour comprendre ce qu'on compte -- il est utile
de savoir \emph{fabriquer} les objets. Les algorithmes ci-dessous produisent
ce que les formules comptent.

\propriete[La relation de Pascal comme algorithme]{
Pour $1 \leq k \leq n-1$,
\[
  \binom{n}{k} = \binom{n-1}{k-1} + \binom{n-1}{k},
  \qquad
  \binom{n}{0} = \binom{n}{n} = 1 .
\]
Chaque ligne se deduit donc de la precedente par des \textbf{additions
seulement} : aucune factorielle, aucune division. Entree : $n$. Sortie : la
liste des $n+1$ coefficients de la ligne $n$.
}

\begin{python}
def ligne_de_pascal(n):
    """Liste des coefficients binomiaux de la ligne n, par additions."""
    if n < 0:
        raise ValueError("n doit etre positif")
    ligne = [1]
    for _ in range(n):
        ligne = [1] + [a + b for a, b in zip(ligne, ligne[1:])] + [1]
    return ligne

def permutations(elements):
    """Toutes les permutations d'une liste, par insertions successives."""
    if not elements:
        return [[]]
    resultat = []
    for petite in permutations(elements[1:]):
        for position in range(len(petite) + 1):
            resultat.append(petite[:position] + [elements[0]] + petite[position:])
    return resultat

def parties(elements, taille):
    """Toutes les parties d'une taille donnee, dans l'ordre des indices."""
    if taille < 0 or taille > len(elements):
        return []
    if taille == 0:
        return [[]]
    avec = [[elements[0]] + reste for reste in parties(elements[1:], taille - 1)]
    return avec + parties(elements[1:], taille)
\end{python}

\exemple{
\code{ligne\_de\_pascal(6)} renvoie $[1, 6, 15, 20, 15, 6, 1]$ : on y lit
$\binom{6}{2} = 15$ sans jamais calculer $6!$.

\code{permutations([1, 2, 3, 4])} produit $24 = 4!$ listes, toutes distinctes.

\code{parties([1, 2, 3, 4, 5], 2)} produit $10 = \binom{5}{2}$ paires.
}

\propriete[Travail demande]{%
\begin{enumerate}
  \item Verifier que la somme des coefficients de la ligne $n$ vaut $2^n$, et
        expliquer pourquoi.
  \item Verifier que chaque ligne est symetrique.
  \item Verifier que \code{parties} produit bien $\binom{n}{k}$ listes, sans
        repetition.
  \item Combien de permutations pour $n = 10$ ? Est-il raisonnable de toutes
        les engendrer ?
\end{enumerate}
}

\exemple{
\textbf{Reponses.} 1. La somme vaut $2^n$ : c'est le nombre total de parties
d'un ensemble a $n$ elements, chaque partie ayant une taille $k$ et une seule.
4. $10! = 3\,628\,800$. Les engendrer toutes est possible mais couteux ; pour
$n = 20$ ce serait $2{,}4\times10^{18}$ listes, hors d'atteinte. Denombrer
n'est pas enumerer -- c'est justement pourquoi on denombre.
}

% BEGIN-VERIFY
% from math import comb, factorial
%
% def ligne_de_pascal(n):
%     if n < 0:
%         raise ValueError("n doit etre positif")
%     ligne = [1]
%     for _ in range(n):
%         ligne = [1] + [a + b for a, b in zip(ligne, ligne[1:])] + [1]
%     return ligne
%
% def permutations(elements):
%     if not elements:
%         return [[]]
%     resultat = []
%     for petite in permutations(elements[1:]):
%         for position in range(len(petite) + 1):
%             resultat.append(petite[:position] + [elements[0]] + petite[position:])
%     return resultat
%
% def parties(elements, taille):
%     if taille < 0 or taille > len(elements):
%         return []
%     if taille == 0:
%         return [[]]
%     avec = [[elements[0]] + reste for reste in parties(elements[1:], taille - 1)]
%     return avec + parties(elements[1:], taille)
%
% assert ligne_de_pascal(6) == [1, 6, 15, 20, 15, 6, 1]
% # L'algorithme redonne exactement les coefficients binomiaux.
% for n in range(0, 15):
%     ligne = ligne_de_pascal(n)
%     assert ligne == [comb(n, k) for k in range(n + 1)]
%     assert len(ligne) == n + 1
%     assert sum(ligne) == 2 ** n
%     assert ligne == ligne[::-1]
%     # ... et la relation de Pascal elle-meme.
%     if n >= 1:
%         precedente = ligne_de_pascal(n - 1)
%         assert all(
%             ligne[k] == precedente[k - 1] + precedente[k] for k in range(1, n)
%         )
%
% p4 = permutations([1, 2, 3, 4])
% assert len(p4) == factorial(4) == 24
% assert len({tuple(x) for x in p4}) == 24
% assert all(sorted(x) == [1, 2, 3, 4] for x in p4)
% for n in range(0, 6):
%     assert len(permutations(list(range(n)))) == factorial(n)
%
% paires = parties([1, 2, 3, 4, 5], 2)
% assert len(paires) == comb(5, 2) == 10
% assert len({tuple(x) for x in paires}) == 10
% assert all(len(x) == 2 for x in paires)
% for n in range(0, 8):
%     for k in range(0, n + 1):
%         engendrees = parties(list(range(n)), k)
%         assert len(engendrees) == comb(n, k), (n, k)
%         assert len({tuple(x) for x in engendrees}) == comb(n, k)
% # Une taille impossible ne produit rien.
% assert parties([1, 2], 3) == []
% assert parties([1, 2], -1) == []
%
% assert factorial(10) == 3628800
% try:
%     ligne_de_pascal(-1)
% except ValueError:
%     pass
% else:
%     raise AssertionError("n negatif aurait du etre refuse")
% END-VERIFY

% ============================================================
% STRATE 2 --- Erreurs frequentes
% ============================================================

\erreurFrequente{
\textbf{Calculer $\binom{n}{k}$ par $\dfrac{n!}{k!(n-k)!}$ dans un
programme.} Pour $n = 100$, $n!$ compte $158$ chiffres alors que
$\binom{100}{50}$ n'en compte que $30$. La relation de Pascal n'utilise que
des additions, et ne fabrique jamais de nombres inutilement grands.
}

\erreurFrequente{
\textbf{Confondre enumerer et denombrer.} Engendrer les $\binom{n}{k}$ parties
permet de \emph{verifier} un denombrement sur de petits cas. Cela ne le
remplace pas : le nombre de parties croit trop vite pour qu'on puisse les
lister au-dela de quelques dizaines d'elements.
}
